\clearpage
\section{Dereferencing through arbitrary atom trees}
\label{sec:copying-dereference}
Figure~19(b) of the paper (PDF p.~38) now has a concrete local dereference
rule for an arbitrary atom contraction tree with $h+1$ inputs, for every
$h\geq0$. One input belongs to the instance being read; the other $h$ inputs
remain available to other references. Both rule sides are actual finite typed
hypernets, not a heap lookup substituted for the paper's graph rule.

The computation keeps the store and the atom tree, weakens the consumed
instance, and inserts $D^{\star}_{1,2}$ between the store and its contents.
Search then requests copying at that term distributor. It does not yet return
the stored value. Exact store resolution, preservation of other reference
inputs, renaming, forward padding, safety, and a five-transition local
search/read execution are proved.

The complete allocation/application/read execution of Proposition~12.2 is
still pending. In particular, combining the local allocation and read rules
with arbitrary focused contexts, proving global rule matching and determinism,
and supplying the ensuing copy transitions remain necessary. The local read
execution below does not claim to be a run of the original translated entire
program. The terms and all Figure~16 translation clauses remain unchanged.

\Needspace{12\baselineskip}
\subsection{Finite wiring frames with exposed reference inputs}

\termdef{copying-read-interface}
The common input list is one term wire followed by $h$ atom wires; the output
is one term wire. A named tree hole expects $h+1$ atom inputs and one atom
output. A second, differently named hole expects two term inputs and one term
output. The first will receive the preserved atom tree, the second the new
term distributor. The redex edge tags are reader, instance, tree and store;
the reduct edge tags are discard, instance, tree, store and share. These are
finite enumeration types, not semantic operations.

\lean{Copying/DereferenceFrames}{9}{9}{Dereference.inputs}

\lean{Copying/DereferenceFrames}{11}{11}{Dereference.treeHole}

\lean{Copying/DereferenceFrames}{12}{12}{Dereference.shareHole}

\lean{Copying/DereferenceFrames}{14}{16}{Dereference.BeforeEdge}

\lean{Copying/DereferenceFrames}{17}{19}{Dereference.AfterEdge}

The boundary-index helper embeds an index of $h+1$ into the length-indexed
input interface without changing its numerical value. The other-reference
list consists of input positions $i+1$, in order, for $0\leq i<h$.
Its no-duplicates proof composes injectivity of the input constructor with
injectivity of successor. Membership in that list is equivalent to being
one of these successor-indexed inputs, by the finite-range map definition.

\lean{Copying/DereferenceFrames}{21}{22}{Dereference.inputPort}

\lean{Copying/DereferenceFrames}{24}{25}{Dereference.otherInputs}

\lean{Copying/DereferenceFrames}{27}{28}{Dereference.otherInputs_nodup}

\lean{Copying/DereferenceFrames}{30}{32}{Dereference.mem_otherInputs}

In the redex frame, the reader consumes term input zero and produces the
instance's term input. The instance produces the first atom-tree input; the
remaining atom-tree inputs are exactly the external successor positions.
The tree output enters the store, whose term output is the sole external
output. The frame has three internal vertices and four edges. Their label,
source and target functions explicitly encode this wiring.

The net constructor proves source/target lists duplicate-free, uniqueness of
the incident edge at each side of a vertex, coverage of all non-boundary
positions, and equality of incidence types with each primitive signature.
Finite edge case analysis proves uniqueness. Coverage splits internal
positions into their three cases and input positions into zero or successor.
Typing of the variable-length atom list follows from its constant atom labels.
These are proved fields of the actual graph, not external well-formedness
assumptions.

\lean{Copying/DereferenceFrames}{34}{38}{Dereference.beforeLabel}

\lean{Copying/DereferenceFrames}{40}{44}{Dereference.beforeSource}

\lean{Copying/DereferenceFrames}{46}{50}{Dereference.beforeTarget}

\lean{Copying/DereferenceFrames}{53}{104}{Dereference.beforeFrame}

The reduct replaces the reader by term weakening feeding the old instance.
The instance, atom tree and store remain connected in the same order. The
store's term output and external term input zero feed the two inputs of the
new share hole, in that order. Its output is the external term output. This
frame has four internal vertices and five edges. Its typing and incidence
proofs repeat the explicit finite checks, with the store-result wire and
share hole as the additional cases. The other reference inputs are still
attached to the old atom-tree positions.

\lean{Copying/DereferenceFrames}{106}{111}{Dereference.afterLabel}

\lean{Copying/DereferenceFrames}{113}{118}{Dereference.afterSource}

\lean{Copying/DereferenceFrames}{120}{125}{Dereference.afterTarget}

\lean{Copying/DereferenceFrames}{128}{185}{Dereference.afterFrame}

\Needspace{12\baselineskip}
\subsection{Filling the frames with the actual trees}

\termdef{copying-dereference-rule}
Given a contraction tree $T$ of atom type with $h+1$ inputs, the redex filling
puts its actual net in the tree slot and a primitive corolla at each other
edge. The reduct filling puts the same $T$ in its tree slot, the original
canonical tree $D^{\star}_{1,2}$ in the share slot, and the designated primitive
corollas at the other edges. Simultaneous expansion removes the holes and
connects these components by the proved frame incidences. Relabelling primitives
into allowance $d$ yields both copying hypernets. In particular, no term
weakening or contraction in the canonical tree is omitted.

The raw local rule consists of these two constructed graphs, for any $T$ and
$d$. Its inversion theorem recovers a tree witness and equalities with the two
sides by elimination of the rule constructor; the converse substitutes those
equalities and constructs the rule. Unlike allocation, dereferencing does not
replace or evaluate its external stored contents, so this local fragment has
no stable-value parameter. The subsequent copying operation has its own
copyability obligation.

\lean{Copying/Dereference}{10}{17}{Dereference.beforeFill}

\lean{Copying/Dereference}{20}{28}{Dereference.afterFill}

\lean{Copying/Dereference}{31}{33}{Dereference.before}

\lean{Copying/Dereference}{36}{38}{Dereference.after}

\lean{Copying/Dereference}{42}{43}{Dereference.Rule}

\lean{Copying/Dereference}{45}{49}{Dereference.rule_iff}

The instance and store occurrences on each side have explicit expanded
edge tags. The tree injections retain every original tree edge under its
component tag. Their label equation is literal preservation of the old tree
label, lifted to the copying alphabet. One allowance raise of either rule
side equals constructing that same side at allowance $d+1$; this follows by
cases on the allowance, since every edge is primitive. The rule therefore
survives a raise by rewriting these two equalities.

The first search step on the reduct meets the second contraction of
$D^{\star}_{1,2}$ at external term input zero. The exact source membership and
term-contraction label give a rewrite token at that same vertex. This is the
required copy request.

\lean{Copying/Dereference}{51}{51}{Dereference.beforeInstance}

\lean{Copying/Dereference}{52}{52}{Dereference.afterInstance}

\lean{Copying/Dereference}{53}{53}{Dereference.beforeStore}

\lean{Copying/Dereference}{54}{54}{Dereference.afterStore}

\lean{Copying/Dereference}{55}{55}{Dereference.beforeTree}

\lean{Copying/Dereference}{56}{56}{Dereference.afterTree}

\lean{Copying/Dereference}{58}{60}{Dereference.tree_labels}

\lean{Copying/Dereference}{63}{64}{Dereference.before_raise}

\lean{Copying/Dereference}{66}{67}{Dereference.after_raise}

\lean{Copying/Dereference}{69}{72}{Dereference.rule_raise}

\lean{Copying/Dereference}{75}{79}{Dereference.requests_copy}

\Needspace{12\baselineskip}
\subsection{Preservation of resolved store identity and trace length}

\termdef{copying-read-resolution}
Each side has the actual expansion attachment map from tree vertices into
the surrounding fragment. If two tree edges are adjacent, map their shared
vertex through this attachment. Mapped target/source membership proves
adjacency in the fragment. The tree's weakening-or-contraction label condition
makes this an atom hop. The argument applies separately to the redex and
reduct with identical original tree edges.

The instance enters whichever tree edge consumes input zero: its target and
that inserted input are the same frame vertex. At the other end, the unique
tree root targets the store's input. Root-output membership supplies the shared
vertex, and the root's permitted label supplies the atom-hop condition. Thus
the two boundary hops are proved from the actual graph construction.

\lean{Copying/DereferenceResolution}{9}{10}{Dereference.beforeTreeVertex}

\lean{Copying/DereferenceResolution}{12}{13}{Dereference.afterTreeVertex}

\lean{Copying/DereferenceResolution}{15}{22}{Dereference.tree_hop_before}

\lean{Copying/DereferenceResolution}{24}{31}{Dereference.tree_hop_after}

\lean{Copying/DereferenceResolution}{34}{39}{Dereference.enters_before}

\lean{Copying/DereferenceResolution}{41}{46}{Dereference.enters_after}

\lean{Copying/DereferenceResolution}{48}{55}{Dereference.root_before}

\lean{Copying/DereferenceResolution}{57}{64}{Dereference.root_after}

Appending one relation step to an $n$-step trace gives an $(n+1)$-step
trace, by induction on the original trace. The arbitrary-tree theorem supplies
an edge consuming input zero and an $n$-step trace from that edge to its root.
Map this trace into each rule side, append the root-to-store hop, and prepend
the instance-to-tree hop. Both sides resolve to their designated store in
$n+2$ atom hops.

This proves existence without choosing a canonical tree presentation. Given
any proposed successful resolution at cost $m$ on either side, deterministic
atom tracing and terminal-store uniqueness identify $m$ with the common cost
above. Consequently the two sides admit successful resolution at exactly the
same counts. The same uniqueness theorem excludes every other endpoint store.
These counts measure atom-trace witnesses, not UAM execution steps. This does
not yet prove uniqueness of all possible context matches in a whole machine
state, nor identify non-isomorphic tree presentations as graphs.

\lean{Copying/DereferenceResolution}{67}{71}{Dereference.trace_snoc}

\lean{Copying/DereferenceResolution}{74}{81}{Dereference.resolves_both}

\lean{Copying/DereferenceResolution}{84}{90}{Dereference.resolution_count_iff}

\lean{Copying/DereferenceResolution}{93}{96}{Dereference.store_unique_before}

\lean{Copying/DereferenceResolution}{99}{102}{Dereference.store_unique_after}

\Needspace{12\baselineskip}
\subsection{The local execution and its pending copy}

The reader's argument vertex is the first internal frame vertex. Entering
the actual reader edge searches this vertex; entering the actual instance
returns backtrack there. Returning through the unary reader then changes to
rewrite mode at external term input zero. All three search witnesses carry
their exact edge labels and ordered incidence positions.

\lean{Copying/DereferenceExecution}{9}{9}{Dereference.argument}

\lean{Copying/DereferenceExecution}{11}{15}{Dereference.search_argument}

\lean{Copying/DereferenceExecution}{17}{20}{Dereference.instance_returns}

\lean{Copying/DereferenceExecution}{22}{24}{Dereference.ready}

\termdef{copying-read-subsystem}
A state pairs an actual graph at the local read interface with a token
coordinate in that graph. The subsystem relation consists of intrinsic search
and the raw read computation from rewrite mode on the redex to search mode on
the reduct. Successful finality remains backtrack at an external input.
Such states have no search successor by the input obstruction and no read
successor because the read constructor requires rewrite mode. This proves
final normality for the subsystem machine.

Concatenate the three search transitions, the single read transition, and the
proved term-contraction search transition. The result is a five-step execution
ending in rewrite mode on the reduct. Its finality predicate is false, because
rewrite and backtrack modes differ. No claim of combined-machine determinism
or value return is hidden in this execution theorem.

\lean{Copying/DereferenceExecution}{27}{27}{Dereference.State}

\lean{Copying/DereferenceExecution}{30}{33}{Dereference.Step}

\lean{Copying/DereferenceExecution}{35}{44}{Dereference.final_normal}

\lean{Copying/DereferenceExecution}{46}{49}{Dereference.machine}

\lean{Copying/DereferenceExecution}{52}{57}{Dereference.read_run}

\lean{Copying/DereferenceExecution}{60}{62}{Dereference.read_pending_copy}

Two further incidence equalities check the intended read action: the old
instance source is now the target of the discard weakening; and the new
term-contraction source list consists of the store output followed by the
external result input. For each other reference $i<h$, the tree's input $i+1$
attaches to external atom position $i+1$ on both sides. Unfold the expansion's
input attachment and the ordered finite-range list to prove these equations.
Thus reading neither reconnects another reference to a different input nor
consumes it.

\lean{Copying/DereferenceExecution}{65}{66}{Dereference.instance_discarded}

\lean{Copying/DereferenceExecution}{69}{70}{Dereference.shared_store_result}

\lean{Copying/DereferenceExecution}{73}{76}{Dereference.before_other_input}

\lean{Copying/DereferenceExecution}{78}{81}{Dereference.after_other_input}

\Needspace{12\baselineskip}
\subsection{Carrier renaming, safety, and retained contents}

An equality-labelled isomorphism of primitive nets lifts to a recursive
copying isomorphism at every allowance. Keep its carrier bijections and ordered
incidence equations; label equality followed by reflexivity of recursive label
equivalence proves the label field. To rename a matched tree, assemble its net
isomorphism at the tree component and identity isomorphisms everywhere else.
Simultaneous-expansion congruence supplies the redex and reduct isomorphisms,
and the lifting construction promotes them to allowance $d$. Both maps send
the designated store occurrence to the designated store occurrence by
definition. Iterating the raw one-step padding theorem proves forward padding
of local rules by any finite allowance increment.

\lean{Copying/DereferenceTransport}{9}{19}{Dereference.liftIso}

\lean{Copying/DereferenceTransport}{21}{28}{Dereference.beforeBaseIso}

\lean{Copying/DereferenceTransport}{30}{38}{Dereference.afterBaseIso}

\lean{Copying/DereferenceTransport}{41}{42}{Dereference.before_rename}

\lean{Copying/DereferenceTransport}{44}{45}{Dereference.after_rename}

\lean{Copying/DereferenceTransport}{48}{50}{Dereference.rename_store}

\lean{Copying/DereferenceTransport}{53}{57}{Dereference.rule_pad}

A primitive label satisfies a recursively lifted predicate precisely when
the primitive itself satisfies it, by allowance cases. Hence the redex has
any global predicate true of reader, instance, store, atom weakening and atom
contraction. The reduct has any global predicate true of instance, store and
both sorts of weakening and contraction. Case analysis on the expanded edge
component proves these claims; tree components use the contraction-tree label
condition or the canonical-tree label definition. Specialise to exclusion of
focus labels and hole labels to prove both rule sides focus-free and hole-free.

\lean{Copying/DereferenceSafety}{9}{10}{Dereference.label_primitive}

\lean{Copying/DereferenceSafety}{12}{23}{Dereference.before_labels}

\lean{Copying/DereferenceSafety}{25}{41}{Dereference.after_labels}

\lean{Copying/DereferenceSafety}{43}{49}{Dereference.focusFree_sides}

\lean{Copying/DereferenceSafety}{51}{57}{Dereference.holeFree_sides}

Appending any graph $G:\star\Rightarrow O$ to the reduct gives the read
fragment with actual contents. Each edge in $G$ retains its label through the
right serial injection. The output of the new term distributor is attached
exactly to the old input of $G$, by the actual serial incidence equation.
No assumption that $G$ is acyclic or already evaluated is used here. The
known local resolution trace persists in the left component of serial assembly;
comparing it with any successful trace in the complete graph forces that
trace to end at the retained store. These are contents and lookup laws, not
a completed context-closure theorem for read computations.

\lean{Copying/DereferenceSafety}{60}{61}{Dereference.withContent}

\lean{Copying/DereferenceSafety}{63}{65}{Dereference.content_label}

\lean{Copying/DereferenceSafety}{68}{70}{Dereference.content_attachment}

\lean{Copying/DereferenceSafety}{73}{77}{Dereference.content_store_unique}

\Needspace{12\baselineskip}
\subsection{Canonical and noncanonical regression instances}

For the canonical atom tree with one input, its root is contraction edge
one, as witnessed by its output incidence and uniqueness of the output
producer. The instance enters this same edge, so the empty internal tree
trace plus the two boundary hops gives resolution in exactly two hops.
Resolution-count preservation gives the same result after reading.

For a noncanonical example, use the weakening-free binary contraction tree.
Its sole edge is both entry and root, giving two-hop resolution with one
additional exposed reference. The other-reference theorem keeps that reference
at external atom position one, and the general execution theorem supplies the
five-step read trace for this noncanonical match.

\lean{Copying/DereferenceExamples}{10}{17}{Dereference.canonical_resolves_before}

\lean{Copying/DereferenceExamples}{19}{23}{Dereference.canonical_resolves_after}

\lean{Copying/DereferenceExamples}{26}{35}{Dereference.binary_resolves_before}

\lean{Copying/DereferenceExamples}{38}{40}{Dereference.binary_other_reference}

\lean{Copying/DereferenceExamples}{43}{45}{Dereference.binary_read_run}

Finally attach the original empty-environment translation of literal one
to the canonical reduct. Its named retained value edge still has the constant
one label, and its source equals the term distributor's output target list.
The canonical term distributor has all three expected edges: term weakening
and both term contractions. These equalities are checked on the actual graph
construction. They establish preserved contents and the real copy handoff;
they do not supply the pending copying execution or a complete execution of
$!(\mathtt{ref}\ 1)$ from its term translation.

\lean{Copying/DereferenceExamples}{48}{48}{Dereference.storedOne}

\lean{Copying/DereferenceExamples}{50}{50}{Dereference.storedOneValue}

\lean{Copying/DereferenceExamples}{52}{52}{Dereference.stored_one_label}

\lean{Copying/DereferenceExamples}{54}{55}{Dereference.stored_one_attachment}

\lean{Copying/DereferenceExamples}{58}{62}{Dereference.canonical_copy_edges}

Validation checks all 311 build jobs. The three axiom audits check six
abstract, 477 graph and 489 copying declarations and use only standard Lean
logical axioms; no unfinished proofs or project-specific axioms occur.
Estimated completion of Proposition~12.2 is about 60\% by mathematical effort,
with broad uncertainty of 45--70\% (26 October 2026). This judgement concerns
the remaining mathematics. Copying, integration and matching of the complete
machine, and the template robustness proofs remain substantial obligations.

